\chapter{The Real Numbers}\label{ch:reals}
\chapterquote{Numbers are free creations of the human mind; they serve as a
means of apprehending more easily and more sharply the difference of things.}
{Richard Dedekind}{Was sind und was sollen die Zahlen?}

Analysis begins with the real numbers. In this chapter we list the axioms
that characterize $\R$ as a complete ordered field and extract the first
consequences that will be used throughout the book: the Archimedean property
and the density of the rationals.\footnote{Both go back to antiquity: the
Archimedean property appears in Euclid as Definition V.4, and it is used by
Eudoxus in the theory of proportion.}

The aim is not completeness of exposition but economy: everything proved
here is used again in later chapters.

% ------------------------------------------------------------------
\section{Ordered fields}\label{sec:ordered-fields}

\subsection{Fields}
We begin with the algebraic structure.

\begin{definition}[Field]\label{def:field}
A \emph{field} is a set $F$ with two binary operations $+$ and $\cdot$ such
that $(F,+)$ is an abelian group with identity $0$, $(F\setminus\{0\},\cdot)$
is an abelian group with identity $1$, and multiplication distributes over
addition.
\end{definition}

\begin{example}[The rationals]\label{ex:rationals}
$\Q$ with its usual operations is a field. So is $\Z/p\Z$ for every prime $p$;
we will see in a moment that the latter cannot be ordered.
\end{example}

\subsection{Order}
The real numbers carry more than an algebraic structure.

\begin{definition}[Ordered field]\label{def:ordered-field}
An \emph{ordered field} is a field $F$ together with a subset $P\subseteq F$
(the \emph{positive} elements)\footnote{Equivalently, one may specify a total
order $<$ compatible with $+$ and $\cdot$.} closed under $+$ and $\cdot$, such that for
every $x\in F$ exactly one of $x\in P$, $x=0$, $-x\in P$ holds. We write
$x<y$ when $y-x\in P$.
\end{definition}

\begin{proposition}\label{prop:squares}
In an ordered field, $x^2>0$ for every $x\neq 0$. In particular $1>0$.
\end{proposition}

\begin{proof}
This is \cref{ex:squares}.
\end{proof}

\begin{corollary}\label{cor:no-order-Zp}
The field $\Z/p\Z$ admits no ordering.
\end{corollary}

\begin{proof}
If it did, $1>0$ by \cref{prop:squares}, hence
$0 = 1 + 1 + \cdots + 1 > 0$ ($p$ summands), a contradiction.
\end{proof}

\begin{exercises}
\begin{exercise}
Prove that the additive and multiplicative identities of a field are unique.
\end{exercise}

\begin{exercise}\label{ex:squares}
Prove \cref{prop:squares}: in an ordered field $x^2>0$ for all $x\neq 0$.
\fromtext{prop:squares, sec:completeness}
\end{exercise}

\begin{exercise}\label{ex:abs-value}
Define $|x| = \max(x,-x)$ in an ordered field. Prove the triangle inequality
$|x+y|\le|x|+|y|$.
\fromex{ex:reverse-triangle}
\end{exercise}

\begin{exercise}[Complex numbers]
Show that $\mathbb{C}$ admits no ordering. (Use \cref{prop:squares}.)
\end{exercise}
\end{exercises}

% ------------------------------------------------------------------
\section{Completeness}\label{sec:completeness}

\subsection{Suprema}
\begin{definition}[Supremum]\label{def:sup}
Let $S$ be a subset of an ordered field $F$. An element $s\in F$ is the
\emph{supremum} of $S$, written $s=\sup S$, if $s$ is an upper bound of $S$
and $s\le t$ for every upper bound $t$ of $S$.
\end{definition}

\begin{principle}[Completeness axiom]\label{pr:completeness}
Every nonempty subset of $\R$ that is bounded above has a supremum in $\R$.
\end{principle}

\begin{remark}
\Cref{pr:completeness} is what separates $\R$ from $\Q$: the set
$\{q\in\Q : q^2<2\}$ is bounded above in $\Q$ but has no supremum there.
\end{remark}

\subsection{Consequences}
\begin{theorem}[Archimedean property]\label{thm:archimedean}
For every $x\in\R$ there is $n\in\N$ with $n>x$. Equivalently,
\begin{equation}\label{eq:archimedean}
  \inf_{n\ge 1}\frac1n = 0.
\end{equation}
\end{theorem}

\begin{proof}
If not, $\N$ is bounded above; let $s=\sup\N$ by \cref{pr:completeness}.
Then $s-1$ is not an upper bound, so $n>s-1$ for some $n\in\N$, whence
$n+1>s$ with $n+1\in\N$, a contradiction.
\end{proof}

\begin{corollary}[Density of $\Q$]\label{cor:density}
Between any two real numbers $a<b$ there is a rational number.
\end{corollary}

\begin{proof}
By \eqref{eq:archimedean} choose $n$ with $1/n<b-a$; now apply
\cref{ex:density}.
\end{proof}

\begin{example}[An irrational supremum]
The set $S=\{q\in\Q: q^2<2\}$ has $\sup S=\sqrt2$ in $\R$.
\end{example}

\begin{exercises}
\begin{exercise}\label{ex:density}
Let $a<b$ and $1/n<b-a$. Show that some $m/n$ lies in $(a,b)$.
\fromtext{cor:density}
\end{exercise}

\begin{exercise}\label{ex:reverse-triangle}
Prove the reverse triangle inequality $\bigl||x|-|y|\bigr|\le|x-y|$.
(Use \cref{ex:abs-value}.)
\end{exercise}

\begin{exercise*}[Dedekind cuts]\label{ex:cuts}
Construct $\R$ from $\Q$ by Dedekind cuts, and verify
\cref{pr:completeness} for your construction.
\end{exercise*}
\end{exercises}

% ------------------------------------------------------------------
\begin{problems}
\begin{problem}[Uniqueness of $\R$]\label{prob:unique}
Show that any two complete ordered fields are isomorphic as ordered fields.
\begin{parts}
  \item Show that each contains a unique copy of $\Q$.
  \item Extend the identity of $\Q$ using suprema, and check that it is an
        order-preserving field isomorphism.
\end{parts}
\hint{Density of $\Q$, \cref{cor:density}, is used in both directions.}
\end{problem}

\begin{problem*}\label{prob:non-arch}
Construct an ordered field that is not Archimedean.
\hint{Order the rational functions $\R(t)$ by the sign of the leading
coefficient.}
\end{problem*}
\end{problems}
